% Ported verbatim from paper/top3_joint_moment_bound.tex lines 1192-1227.
% Only cross-reference wording was adapted; no mathematics was altered.
For $D<1/20$ define the explicitly nonnegative quadratic expression
\begin{align}\label{Q8}
 Q={}&\frac{13}{200}\left(\sqrt V-\frac{937}{273}q\sqrt e\right)^2
       +\frac{19}{50}\left(\sqrt J-\frac{2811}{2128}q\sqrt e\right)^2\notag\\
 &+\left(\frac{183}{200}-\frac{725029049}{278853120}q^2\right)e\notag\\
 &+\frac{148831}{181104000}e
       +\frac{42847433}{5229378000000}V+\frac{65082103}{20917512000}J.
\end{align}
The parenthetical coefficient is positive by \eqref{smallbranch} and the
same exact lower bound following \eqref{squares}. The last coefficients are
obtained from those in \eqref{squares}, respectively by adding
$-1/9600$, $-187/1680000+1/375000$, and $211/504000$.
Substitution of \eqref{reserve}--\eqref{paid-expansion} into \eqref{certificate-start} yields the exact identity
\begin{equation}\label{certificate8}
 D=d_1+Q+{\cal R}_0+\ell r_{\rm reg}
          +\frac{K^2}{1125}+r_{\rm tr}+\frac{r_{\rm ct}}9.
\end{equation}
Every remainder is defined from the same $X$, its independent copy and its
complete Gram matrix.
\begin{proposition}\label{conditional8}
Every isotropic law whose radii lie almost surely in a positive bounded
interval of endpoint ratio at most $\sqrt8$ satisfies $D\ge d_1$.
In the branch $D<1/20$, \eqref{certificate8} retains all displayed remainders.
\end{proposition}
\begin{proof}
In the small branch all terms after $d_1$ are nonnegative. In the other
branch $D\ge1/20>d_1$. To verify the last comparison entirely by rationals,
use $\sqrt5<161/72$ and $\sqrt3>433/250$ in \eqref{constants}; they give
\[
 \frac1{20}-d_1>\frac{1025292188593}{804826176000000}>0.
\]
No division by $B=0$ occurs in the small branch, and the other branch
requires none. The radial hypothesis will be verified for a maximizing law in
Section~\ref{sec-conclusion}, rather than imposed on every isotropic law.
\end{proof}
